\section{Conclusion}
\label{sec:conclusion}

We present a systematic computational study demonstrating that trace impurities in lithium batteries are not universally harmful. Through physics-based simulations and analysis of published experimental data, we show that magnesium at ${\sim}1{,}000$~\ppm improves both Coulombic Efficiency ($+0.0018\%$, $p=0.018$) and capacity retention ($+0.10\%$) by promoting low-resistivity, inorganic-rich SEI. Oxygen-containing species show similar benefits at 200~\ppm. In contrast, iron is detrimental at all concentrations ($-0.33\%$ capacity retention at 1{,}000~\ppm), and water exhibits a threshold effect with a crossover at ${\sim}260$~\ppm. The dominant predictor of impurity benefit is SEI resistivity ($r=-0.891$, $p<0.0001$), yielding a simple design rule: impurities that reduce SEI ionic resistance are beneficial.

These findings open a new perspective on battery manufacturing, where ``impurity engineering'' replaces the pursuit of ultra-high purity. Controlled introduction of Mg or oxygen-containing species at 200--10{,}000~\ppm could reduce precursor costs while improving cell performance. Future work should validate these predictions experimentally using Li/Cu half-cells with controlled impurity additions, extend simulations to 500+ cycles, and explore multi-impurity interactions and transferability across cell chemistries.
